\originalchapter{量、单位与向量}\label{ch:units}
\originalsection{测量与模型}
用摄影记录落球, 同一物体在不同时间占据不同位置. 这给出运动学的原始数据: 选定的钟所读的时间 $t$, 与尺所读的位置. 要讨论力, 还需通过惯性或称量给物体赋予质量 $m$. 长度、时间、质量的 SI 单位分别是米 (m)、秒 (s)、千克 (kg). 质量反映改变速度的困难程度; 重量则是引力, 两者不是同一种量.

质点不是尺寸为零的真实物体, 而是忽略尺寸和内部形变后的模型. 讨论汽车跨城行驶可把车作为质点; 讨论轮胎滚动就必须保留半径与转动. 刚体模型要求任意两物质点的距离保持不变. 弹簧、绳与碰撞的理想化条件会在使用处另给, 不能仅从一幅图判断它们是否无质量或无摩擦.

量纲 $[X]$ 描述量 $X$ 对基本单位的依赖. 速度量纲为 $LT^{-1}$, 加速度为 $LT^{-2}$, 力为 $MLT^{-2}$. 等式两边及加减的每一项必须同量纲; 三角、指数函数的自变量必须无量纲. 弧度是弧长与半径之比, 因而无量纲, 但写 rad 可提示角度的物理意义. 写 $\sin(\omega t)$ 时, $\omega$ 必须有 $T^{-1}$ 的量纲.

\begin{example}
仅知单摆长度 $L$, 重力加速度 $g$, 摆球质量 $m$ 与初始角 $\theta_0$, 估计周期的依赖.
\end{example}
\begin{solution}
设 $T=Cm^aL^bg^c$. 比较量纲得到 $a=0$, $b+c=0$, $-2c=1$, 所以 $T=\sqrt{L/g}\,f(\theta_0)$. 无量纲角可以进入未知函数, 量纲分析既不能给出系数 $2\pi$, 也不能证明周期与振幅无关. 小角度下的具体函数来自第\ref{ch:osc}章运动方程.
\end{solution}
估算首先声明输入量的数量级和模型. 例如距离除以时间得到平均速率, 不能把人的身高当作每人的臂展而不说明假设. 数量级估算检验可行性, 不代替需要精确常数或边界条件的解答.

\originalsection{矢量及投影}
位移同时有大小与方向. 在选定的正交坐标系中, 写 $\vct r=x\uv x+y\uv y+z\uv z$, 其中 $\uv x,\uv y,\uv z$ 是长度为 1 的方向向量, 而 $x,y,z$ 是可正可负的分量. 向量本身不因坐标系的转动而改变; 分量会改变. $|\vct r|=\sqrt{x^2+y^2+z^2}$ 是坐标无关的长度. 向量加法把连续两次位移合成为总位移, 减法给出从一位置到另一位置的位移.

点积 $\vct a\cdot\vct b=ab\cos\theta$ 给出一个标量, 正交坐标下等于 $a_xb_x+a_yb_y+a_zb_z$. 它将一个向量在另一个方向上投影, 因而适合表达力对位移的功. 叉积 $\vct a\times\vct b$ 垂直于二者平面, 大小 $ab\sin\theta$, 方向按右手定则: 手指由 $\vct a$ 转向 $\vct b$, 拇指指向叉积. 调换次序改变符号. 在右手坐标系中, $\uv x\times\uv y=\uv z$, 按循环置换得到其他关系.

\begin{definition}
点 $O$ 到质点的向量 $\vct r$ 与作用于质点的力 $\vct F$ 的叉积 $\vct\tau_O=\vct r\times\vct F$ 称为力关于 $O$ 的力矩. 下标 $O$ 不能省略其参考点含义.
\end{definition}
力矩大小也等于力乘以参考点到力作用线的垂直距离. 力作用线通过参考点时力矩为零, 即使力本身很大. 这解释为何推门时在靠近铰链处用力效果较小. 后文角动量同样需要说明参考点.

\originalsection{导数与积分}
给出位置函数后, $\Delta\vct r/\Delta t$ 表示一段时间的平均速度. 缩短时间段得到导数 $\vct v=\mathrm d\vct r/\mathrm dt$, 它是瞬时速度. 加速度 $\vct a=\mathrm d\vct v/\mathrm dt$ 描述速度向量的变化, 既可能改变速率也可能改变方向. 即使速率恒定, 转弯仍有加速度.

积分把瞬时变化累积起来: $\vct r(t)-\vct r(t_0)=\int_{t_0}^t\vct v(s)\dd s$. 速度图与时间轴之间的有向面积给位移, 速率 $|\vct v|$ 的积分给路程. 回到原地的位移为零, 路程通常非零. 初值明确时才能从加速度唯一还原速度与位置. 后文用 $\dot x$ 表示 $\mathrm dx/\mathrm dt$, $\ddot x$ 表示二阶时间导数; 空间导数则写 $\mathrm dU/\mathrm dx$, 不混用.
\originalsection{单位体系与坐标计算}
经典力学从伽利略的运动测量、牛顿的运动与引力规律逐渐形成. 它将复杂物体抽象成可检验的数学模型, 但并非所有尺度上精确: 原子尺度常需量子理论, 高速运动需相对论. 力学中的三基本量已足够给多数公式定量, 完整 SI 还包括以下七基本单位. 现代 SI 通过固定常数定义, 不再以一个可改变的实物千克样品作为最终标准. 定义以 \href{https://www.bipm.org/en/measurement-units/si-base-units}{BIPM 的 SI 基本单位说明} 为准, 采用当前 SI 定义并说明教材中的历史讨论.
\begin{longtable}{@{}p{.15\textwidth}p{.15\textwidth}p{.62\textwidth}@{}}
\toprule 基本量 & 单位 & 固定常数所确定的尺度 \\\midrule
时间 & s & 铯 133 未扰动基态超精细跃迁频率 $\Delta\nu_{Cs}=9\,192\,631\,770$ Hz. \\
长度 & m & 真空光速 $c=299\,792\,458$ m/s, 配合秒确定米. \\
质量 & kg & 普朗克常数 $h=6.62607015\times10^{-34}$ J s, J s 的量纲为 kg m$^2$/s. \\
电流 & A & 基本电荷 $e=1.602176634\times10^{-19}$ C, C=A s. \\
温度 & K & 玻尔兹曼常数 $k_B=1.380649\times10^{-23}$ J/K. \\
物质的量 & mol & 阿伏伽德罗常数 $N_A=6.02214076\times10^{23}$ mol$^{-1}$, 要说明所数的实体. \\
发光强度 & cd & 频率 $540\times10^{12}$ Hz 的单色辐射光效为 683 lm/W, 配合立体角确定 cd. \\\bottomrule
\end{longtable}
所有以上数值在定义中均为精确值; 实际测量仍有不确定度. 立体角 $\Omega=A/r^2$ 为球面区域面积除以半径平方, 单位 sr, 整球为 $4\pi$ sr. sr 与 rad 均无量纲但角色不同. 万有引力常数 $[G]=L^3M^{-1}T^{-2}$ 则须实验测定, 不是 SI 的精确定义常数.

质量也可在相同校准力下比较加速度: $m_1a_1=m_2a_2$, 故惯性质量比为 $a_2/a_1$. 称量比较重力质量, 二者相同须实验检验. 单位定义和某一次天平的读数是两个层次.

转动坐标轴而不转动物理向量, 若新 $x'$ 轴相对原轴逆时针转 $\beta$, 投影给 $a_{x'}=a_x\cos\beta+a_y\sin\beta$, $a_{y'}=-a_x\sin\beta+a_y\cos\beta$. 主动把向量转 $\beta$ 而保持轴不动则符号反过来. 点积不变是长度与角度不变的结果. 叉积分量为
\begin{equation}
\vct a\times\vct b=(a_yb_z-a_zb_y)\uv x+(a_zb_x-a_xb_z)\uv y+(a_xb_y-a_yb_x)\uv z.
\end{equation}
由分量直接可核验反交换和对加法的分配律. 三重积 $\vct a\times(\vct b\times\vct c)=\vct b(\vct a\cdot\vct c)-\vct c(\vct a\cdot\vct b)$ 常用于惯量推导. 柱坐标在每点有右手基 $\uv r\times\uv\theta=\uv z$, 但 $\uv r,\uv\theta$ 随位置改变.

柱坐标 $(r,\theta,z)$ 用到轴的距离、绕轴方位角和高度定位, $x=r\cos\theta,y=r\sin\theta$, $z$ 保持原高度. $\uv r=(\cos\theta,\sin\theta,0)$, $\uv\theta=(-\sin\theta,\cos\theta,0)$; 它是平面极坐标加一个高度, 不是球坐标的离原点距离. 圆柱体积元 $r\dd r\dd\theta\dd z$ 中的 $r$ 来自小弧长 $r\dd\theta$.

一个完整的 Fermi 估算是学校一天饮水质量. 假设 1000 人、每人两瓶、每瓶 0.5 L, 水密度约 1 kg/L, 分解为人数乘瓶数乘每瓶体积乘密度, 得 $1000\times2\times0.5\times1=10^3$ kg/天. 人数和饮水量各错两倍会使结果相差几倍, 但一般不会把吨级变成克级. 先给输入与边界再相乘, 最后保留相称的有效数字, 才使数量级估算可核查; 这里不是实测学校统计.
